% ============================================================
% CHAPTER 8: WEEKLY BREAKDOWN
% MCDA 5585 / 5586 Major Project Report
% Bhavik Kantilal Bhagat | A00494758 | Saint Mary's University
% ============================================================

\chapter{Weekly Breakdown}
\label{chap:weekly_breakdown}

This chapter provides a phase-by-phase narrative of the internship delivery,
organised around the four major phases that structured the 24-week engagement.
The narrative complements the detailed technical account in
Chapter~\ref{chap:implementation} by focusing on the chronological flow ---
what was done each week, what blocked or accelerated progress, and how the
priorities shifted in response to production events.

% ────────────────────────────────────────────────────────────
\section{Phase 1: Foundation \& CloudWatch Dashboards (Weeks 1--7, Mar~15 -- May~10)}
\label{sec:wb_phase1}
% ────────────────────────────────────────────────────────────

\subsection{Weeks 1--4: Onboarding and First Deliverables (Mar 15 -- Apr 11)}

The engagement opened in part-time mode with approximately 10--20 hours per
week. The first priority was access provisioning: AWS console, JIRA,
Confluence, and CodeCommit access requests were filed on Day~1 and cleared
over the following two weeks. The initial period was used productively for
architecture documentation review on Confluence --- building a working mental
model of the six platforms (\texttt{rpacfs1}, \texttt{ia-bpo},
\texttt{automationhub}, \texttt{meridian}, \texttt{citcoworks}) and the
Lambda, ECS, and EC2 workloads distributed across them.

\begin{table}[ht]
\centering
\caption{Phase 1 --- Weekly Delivery Summary}
\label{tab:wb_phase1}
\begin{tabular}{p{1.2cm}p{4.0cm}p{7.5cm}}
\toprule
\textbf{Week} & \textbf{Dates} & \textbf{Key Deliverables} \\
\midrule
1 & Mar 15--21 & Access provisioning, architecture study, AWS environment orientation \\
2 & Mar 22--28 & IISS-435: RPACFS1 CloudWatch pipeline health dashboard (262 tests) \\
3 & Mar 29--Apr 4 & IISS-435 completion; IISS-455 Lambda dashboard design started \\
4 & Apr 5--11 & IISS-455: Lambda dashboard deployed (51 functions); IISS-461: Confluence inventory page \\
5 & Apr 12--18 & Training week (RPA, CitcoWorks, Meridian, Aexeo KTs); Apr 17 incident observed; IISS-469 started \\
6 & Apr 19--25 & IISS-469 completed; IISS-487 Grafana story created; IISS-482 board set up \\
7 & Apr 26--May 2 & IISS-484 ECS dashboard; IISS-485 EKS dashboard; IISS-483 EC2 gap analysis \\
\bottomrule
\end{tabular}
\end{table}

\textbf{Week 2} produced the first concrete deliverable: the RPACFS1
CloudWatch pipeline health dashboard (IISS-435), built against the
\texttt{rpacfs1-cais-pricing-extract-genai-infrastructure} repository. The
dashboard included Lambda function monitoring, Logs Insights error widgets,
and 262 automated tests.

\textbf{Week 5} (Apr 13--18) was the training week. Four knowledge-transfer
sessions covered the RPA support model (Horace), CitcoWorks and Event Store
(Bhanu Teja Murari), Meridian Kafka architecture (Meridian team), and the
\AE xeo Treasury platform. The the Meridian Incident occurred during this
week, making the observability gap visible firsthand. By week end, IISS-469
was underway and the architectural limitation of the \texttt{SEARCH()}-based
dashboard (the visibility gap for inactive functions) had been documented.

\textbf{Week 7} built the ECS (IISS-484), EKS (IISS-485), and EC2 (IISS-483)
monitoring dashboards. The EC2 gap analysis discovered the
\texttt{BudgetCode=``CFS/RPA''} tag mismatch on \texttt{rpacfs1} instances,
introducing the \texttt{AppName} fallback pattern that became foundational for
the inventory API.

\subsection{Phase 1 Total: 100 hours --- Key Insight}

By the end of Phase~1, the Lambda inventory across all six platforms was complete
and the SEARCH() visibility gap had established the design constraint for the
next phase: CloudWatch Metric Insights SQL rather than subscription filters.

% ────────────────────────────────────────────────────────────
\section{Phase 2: Grafana AMG Intensive Build (Weeks 8--11, May~6 -- May~27)}
\label{sec:wb_phase2}
% ────────────────────────────────────────────────────────────

\begin{table}[ht]
\centering
\caption{Phase 2 --- Weekly Delivery Summary}
\label{tab:wb_phase2}
\begin{tabular}{p{1.2cm}p{4.0cm}p{7.5cm}}
\toprule
\textbf{Week} & \textbf{Dates} & \textbf{Key Deliverables} \\
\midrule
8 & May 3--9 & RPA Support Rotation \#1; IISS-487 elevated to Critical; Grafana scaffolding \\
9 & May 10--16 & May 11: 4 CFN stacks deployed; Grafana workspace live; Container Insights enabled; Grafana upgrade 10.4→12.4; Lambda/ECS/SQS rows \\
10 & May 17--23 & MSK, DocumentDB, OpenSearch, RDS rows; Secrets Manager 64KB limit → S3 migration \\
11 & May 24--30 & Redis row; SNS 3-topic setup; 17 alarms; 8 composite alarms; Grafana alerting \\
\bottomrule
\end{tabular}
\end{table}

\textbf{Week 8} was the first RPA support rotation. IISS-487 was elevated
to Critical priority on May~6. Infrastructure deployment could not begin
until the following Monday due to the support shift, but the week was used
to scaffold the Grafana repository and plan the four-layer architecture.

\textbf{Week 9} (May 10--16) was the single most infrastructure-dense week
of the internship. May~11 saw four CloudFormation stacks deployed, the Grafana
AMG workspace provisioned with its 16-role IAM agent permission model, and
Container Insights enabled on five of the six Meridian clusters. Two days
later, on May~13, the unexpected Grafana 10.4→12.4 upgrade broke all SQL
panels and variable interpolation --- approximately one working day was
consumed migrating queries to the stable Grafana~12.4 model.

By \textbf{Week 11}, the core platform was functional: MSK consumer lag
alarm active, composite alarms configured, and the first six dashboard rows
complete. The detection window for a Meridian Kafka congestion event had
dropped from the 3.5-hour Apr~17 gap to a projected sub-30-second alarm fire.

\subsection{Phase 2 Total: 157 hours --- Key Milestone}

The MSK consumer lag alarm armed itself against live production traffic
on May~11, 2026 --- the inflection point where the observability gap
identified by the Meridian Incident was closed.

% ────────────────────────────────────────────────────────────
\section{Phase 3: Dashboard Completion, Hotfixes \& SQS (Weeks 12--15, May~28 -- Jun~18)}
\label{sec:wb_phase3}
% ────────────────────────────────────────────────────────────

\begin{table}[ht]
\centering
\caption{Phase 3 --- Weekly Delivery Summary}
\label{tab:wb_phase3}
\begin{tabular}{p{1.2cm}p{4.0cm}p{7.5cm}}
\toprule
\textbf{Week} & \textbf{Dates} & \textbf{Key Deliverables} \\
\midrule
12 & May 31--Jun 6 & RPA Support Rotation \#2; IISS-538: MESO autoscale hotfix \\
13 & Jun 7--13 & IISS-706: CAIS Deadline hotfix (7h); EC2 rows; IISS-507 scaffolded (2.5h); IISS-482 SQS (960 tests) \\
14 & Jun 14--20 & Dashboard deep-links; final polish; v728 deployment (140h milestone); IISS-482 closed \\
15 & Jun 21--27 & MSK/EC2 refinements; 6 resource mapper Lambdas; IISS-507 Lambda discoverers (119 functions) \\
\bottomrule
\end{tabular}
\end{table}

\textbf{Week 12} was the second RPA support rotation. The MESO ECS
autoscale issue (IISS-538) was resolved alongside support duties.

\textbf{Week 13} was the most incident-dense week of the internship.
The CAIS Deadline \texttt{SpnegoError} (IISS-706) required 7 hours of
investigation and hotfix work, identifying the module-import credential
caching pattern as the root cause. In parallel, the IISS-507 inventory
API was scaffolded using the Kiro spec-driven workflow (2.5 hours), and
the SQS monitoring feature (IISS-482) was completed with 960 automated tests.

\textbf{Week 14} delivered the final dashboard at version v728, closing
IISS-487 at exactly 140 hours. Deep-link panels, DocumentDB and Redis
restructuring, and the main branch merge completed the primary deliverable.

\subsection{Phase 3 Total: 160 hours --- Key Milestone}

IISS-487 closed at v728 on Jun~18, 2026. The full platform --- 14 rows,
110+ panels, 38 alarms, 4 CloudFormation stacks --- was in production.

% ────────────────────────────────────────────────────────────
\section{Phase 4: Inventory API, X-Ray Research \& Report (Weeks 16--24, Jun~19 -- Aug~31)}
\label{sec:wb_phase4}
% ────────────────────────────────────────────────────────────

\begin{table}[ht]
\centering
\caption{Phase 4 --- Weekly Delivery Summary}
\label{tab:wb_phase4}
\begin{tabular}{p{1.2cm}p{4.0cm}p{7.5cm}}
\toprule
\textbf{Week} & \textbf{Dates} & \textbf{Key Deliverables} \\
\midrule
16 & Jun 28--Jul 4 & IISS-741: MESO stuck tasks (Critical); IISS-507 ECS/EC2 discoverers; log group research \\
17 & Jul 5--11 & IISS-753: VPL selector investigation; IISS-507 EC2 discoverers; X-Ray specs written \\
18 & Jul 12--18 & IISS-487: Log group utilities; 30 saved queries; Grafana Row 17 (logs panel) \\
19 & Jul 19--25 & Log group PR merged; SRE Observability Platform v2.19.9 released \\
20 & Jul 26--Aug 1 & RPA Support Rotation \#4; IISS-838--844: X-Ray tracing specs for 7 services \\
21 & Aug 2--8 & IISS-824: ML-RPA-Status LLM Agent fixes (11 bugs, 194 tests, 99\% coverage) \\
22 & Aug 9--15 & AMG repo cleanup; README alerting docs; report LaTeX chapters drafted \\
23 & Aug 16--22 & IISS-876: VPL GTL booking fix; report results/conclusions chapters \\
24 & Aug 23--31 & Final report compilation; KT sessions with Horace/Kri; Mitacs renewal confirmed \\
\bottomrule
\end{tabular}
\end{table}

\textbf{Week 16} saw the MESO month-end stuck tasks incident (IISS-741,
Critical) requiring manual retrigger of INNOCAP/SAMLMOS fund processes. ECS
discoverers were completed concurrently, confirming the 22-cluster inventory
and implementing the two-stage AppName filter for the shared BudgetCode environments.

\textbf{Week 18--19} added the CloudWatch Logs observability layer to the
Grafana platform: 30 Logs Insights saved queries across 7 platforms, a
log group URL encoding utility, and Grafana Row~17 for log-derived panels.

\textbf{Week 21} addressed the ML-RPA-Status LLM Agent (IISS-824):
11 critical bugs fixed (OOM Lambda, Athena SSE, SSL, idempotent KB sync),
a 1,140-line monolithic function split into 8 modules, and 194 unit tests
written with 99\% coverage.

\textbf{Weeks 22--24} transitioned to report writing, knowledge transfer,
and handover documentation. The Mitacs BSI renewal for
Sep~1 -- Dec~31, 2026 was confirmed, providing continuity for
the X-Ray tracing and log-group discovery work deferred from this engagement.

\subsection{Phase 4 Total: 480 hours --- Key Milestone}

The 562-test inventory API reached 99\% coverage on Jul~14, 2026, completing
the third and final technical deliverable of the internship.

% ────────────────────────────────────────────────────────────
\section{Internship Hours Summary}
\label{sec:wb_hours}
% ────────────────────────────────────────────────────────────

\begin{table}[ht]
\centering
\caption{Internship Hours by Phase}
\label{tab:wb_hours}
\begin{tabular}{p{5.0cm}p{1.2cm}p{5.5cm}p{1.5cm}}
\toprule
\textbf{Phase} & \textbf{Weeks} & \textbf{Primary Deliverables} & \textbf{Hours} \\
\midrule
Phase 1: Foundation \& CloudWatch & 1--7 & IISS-435, IISS-455/461/469, IISS-482--486 scaffolding & 100 \\
Phase 2: Grafana AMG Build & 8--11 & IISS-487 core (4 CFN stacks, Rows 1--6, alarms) & 157 \\
Phase 3: Dashboard Completion \& Hotfixes & 12--15 & IISS-487 v728, IISS-706, IISS-538, IISS-482 SQS & 160 \\
Phase 4: Inventory API, Research \& Report & 16--24 & IISS-507, IISS-741, IISS-824, X-Ray specs, report & 480 \\
\midrule
\textbf{Total} & \textbf{24} & & \textbf{897} \\
\bottomrule
\end{tabular}
\end{table}

The internship ran from Mar~15 to Aug~31, 2026, logging
\textbf{897 hours} across 24 weeks. The remaining 3 hours to reach the
900-hour commitment are covered by the knowledge-transfer and handover
sessions in the final days of Aug and the first days of September.
Complete itemised work logs are provided in the appendix and in the
accompanying \texttt{WorkLogs.xlsx} submission file.
